% Generated draft: requires mathematical and source review.
\documentclass[11pt]{article}
\usepackage[margin=25mm]{geometry}
\usepackage{fontspec}
\setmainfont{Latin Modern Roman}
\usepackage{amsmath,amssymb,mathtools,unicode-math}
\setmathfont{Latin Modern Math}
\usepackage{enumitem,needspace,xurl,hyperref}
\setlength{\parindent}{0pt}
\setlength{\parskip}{5pt}
\setlength{\emergencystretch}{3em}
\newcommand{\N}{\mathbb{N}}
\newcommand{\Z}{\mathbb{Z}}
\newcommand{\Q}{\mathbb{Q}}
\newcommand{\R}{\mathbb{R}}
\newcommand{\C}{\mathbb{C}}
\newcommand{\F}{\mathbb{F}}
\newcommand{\A}{\mathbb{A}}
\newcommand{\PP}{\mathbb{P}}
\newcommand{\E}{\mathbb{E}}
\newcommand{\eps}{\varepsilon}
\newcommand{\dd}{\,\mathrm{d}}
\begin{document}
\section*{UnsolvedMath ID: 2305046}
\textbf{Source problem code:} catalogue entry
\subsection{Definitions and mathematical statement}
Let $\mathbb D=\{z\in\C:|z|<1\}$ and $\mathbb T=\partial\mathbb D$. A function analytic on $\mathbb D$ has an asymptotic value $a$ at $\zeta\in\mathbb T$ if there is a curve $\eta:[0,1)\to\mathbb D$ with $\eta(t)\to\zeta$ and $f(\eta(t))\to a$. The record does not explicitly say whether $a=\infty$ is admitted along with finite values.

For a nonconstant analytic $f$, let $A(f)$ be the set of boundary points admitting such an asymptotic value, with that convention to be fixed. The MacLane class in the source is
\[
 \mathcal A=\{f:\ f\text{ is nonconstant and analytic in }\mathbb D,
                   \ \overline{A(f)}=\mathbb T\}.
\]
An arc tract, as stated in the source, consists of arcs $\gamma_n\subset\mathbb D$ tending to a nondegenerate subarc $\gamma\subset\mathbb T$ such that
\[
 \min_{z\in\gamma_n}|f(z)|\longrightarrow\infty.
\]
The source specifies the arc convergence only as ``in the obvious fashion.'' For compact arcs, Euclidean Hausdorff convergence is one possible precise reading, but endpoint/parametrized or other convergence conditions are not silently identified with it. Compactness, or an infimum convention in place of the displayed minimum, should also be confirmed.

K. F. Barth's existence question is whether there is $f\in\mathcal A$ having an arc tract and a nowhere-vanishing derivative:
\[
 \exists f\in\mathcal A,\qquad
 f'(z)\ne0\text{ for every }z\in\mathbb D,
 \quad f\text{ has an arc tract as above}.
\]
Here ``non-zero derivative'' is read as nowhere zero; reading it merely as $f'\not\equiv0$ would be automatic for a nonconstant analytic function and would give a different question. The asymptotic-value and arc-convergence conventions remain explicit source-review dependencies.
\subsection{Short English statement}
Can a locally univalent analytic function belong to the MacLane class and still grow uniformly to infinity along arcs approaching a genuine boundary arc? This restates the nowhere-zero-derivative interpretation of Barth's question. The supplied triage says it was open in Hayman's 2018 update and no recent resolution was located; no independent current-status review is claimed.
\subsection{Sources}
\begin{itemize}
\item UnsolvedMath permanent ID 2305046, source code catalogue entry, ``Research Problems in Function Theory --- Problem 5.46.'' Registry curation and metadata: Ulam AI and UnsolvedMath Contributors, CC BY 4.0; updated 2026-08-06.
\item Hayman, \emph{Research Problems in Function Theory} (2018 source update), Problem 5.46, attributed to K. F. Barth: \url{https://arxiv.org/abs/1809.07200}. Supplied source-TeX extraction accessed 2026-07-29. The metadata retains the update for status triage and marks status/release and rights review NEEDS\_REVIEW.
\item MacLane is cited in the source for the class and arc-tract question without complete bibliographic entries. No missing title, date or locator is inferred.
\end{itemize}
\end{document}
